\documentclass[11pt,a4paper]{article}
\usepackage[T1]{fontenc}
\usepackage[utf8]{inputenc}
\usepackage[ngerman]{babel}
\usepackage[margin=2.5cm]{geometry}
\usepackage[hidelinks]{hyperref}
\usepackage{array}
\setlength{\parindent}{0pt}
\setlength{\emergencystretch}{2em}
\renewcommand{\arraystretch}{1.25}
\sloppy
\begin{document}
\shorthandoff{"}
\begin{center}
  {\Large\textbf{Webtechnologien -- Übungsblatt 6}}
\end{center}
\section*{Aufgabe 1: Getting Started with Node.js}
\textbf{Aufgabe:} Einen Express-Server starten, die Antwort der Startseite
als HTML kennzeichnen und einen Text aus der URL in die Antwort übernehmen.
Der Text soll sichtbar sein, ohne als HTML ausgeführt zu werden.

\medskip
\textbf{Ergebnis:} Die TODOs in \texttt{server.js} sind direkt ergänzt.
Die unangetastete Vorlage bleibt als \texttt{server\_ORIGINAL.js} erhalten.
\begin{itemize}
  \item \texttt{app.get('/' , ...)} behandelt die Startseite.
    \texttt{res.setHeader('Content-Type', 'text/html')} sagt dem Client,
    dass die Antwort HTML ist. \texttt{status(200)} bedeutet Erfolg;
    \texttt{send} schickt den Text zurück.
  \item \texttt{app.get('/:text', ...)} erkennt einen variablen Teil der
    URL. Bei \texttt{/mytext} steht in \texttt{req.params.text} der Wert
    \texttt{mytext}. Express dekodiert dabei die URL-Zeichen.
  \item \texttt{escapeHtml(s)} ersetzt \texttt{\&}, \texttt{<},
    \texttt{>}, doppelte und einfache Anführungszeichen durch HTML-Entities.
    Das kaufmännische Und wird zuerst ersetzt, damit neu eingefügte
    Entities nicht nochmals verändert werden. So bleibt fremder Text Text
    und kann kein HTML oder JavaScript einschleusen (XSS).
  \item Der sichere Text wird an
    \texttt{This is a simple application receiving } angehängt.
    Auch diese Route antwortet mit Status 200 und HTML.
\end{itemize}

\textbf{Starten:} Im Ordner \texttt{aufgabe1} zuerst einmal
\texttt{npm install} ausführen (installiert Express). Dann \texttt{node server.js} ausführen und
\url{http://localhost:3000/} aufrufen. Mit Strg+C beenden.
Nur einen der Übungsserver gleichzeitig starten, da alle Port 3000 nutzen.

\medskip
\textbf{Test:} Mit Node.js v24 gestartet, Anfragen mit \texttt{curl}; alle Fälle liefern genau diese Antworten (jeweils \texttt{Content-Type: text/html}):
\begin{center}
\begin{tabular}{|>{\raggedright\arraybackslash}p{5.2cm}|>{\raggedright\arraybackslash}p{8.5cm}|}
\hline
\textbf{Fall} & \textbf{Ergebnis} \\ \hline
GET \texttt{/} & 200; genau \texttt{This is a simple application} \\ \hline
GET \texttt{/mytext} & 200; genau \texttt{This is a simple application receiving mytext} \\ \hline
GET \texttt{/\%3Cb\%3Ex} & 200; \texttt{This is a simple application receiving \&lt;b\&gt;x}; enthält \texttt{\&lt;b\&gt;x}, kein \texttt{<b>} \\ \hline
\end{tabular}
\end{center}

\medskip
{\small \textbf{Hilfsmittel:} Claude Code (Claude Opus 5.5) und OpenAI Codex (GPT-6.1-Sol) für Erstellung, Prüfung und Dokumentation.}
\end{document}
